\subsection{17 февраля. Обратный оператор}

% \begin{example}
% Пусть есть нормированные пространства $E_1$, $E_2$, оператор $A$ отображает $E_1$, $E_2$
% \end{example}

\begin{definition}
    Оператор $A: E_1 \to E_2$ называется \textit{обратимым}, если $\forall y \in Im \: A \; \exists!x \in E_1: \; Ax=y.$
    
    $A^{-1}$ называется \textit{обратным оператором}.
\end{definition}

\begin{definition}
    Оператор $R$ называется \textit{правым обратным} к $A$, если $\forall y \in \im \: A \; ARy=y$, т.е. $AR=I_{Im \: A}$

    Оператор $L$~--- левый обратный оператор определяется аналогичным образом. Понятно что $L = R$. Такое введение определяется тем, что в разных ситуациях проще искать разный тип обратного оператора.
\end{definition}

\begin{lemmanote}
    Правый обратный отвечает за существование решения уравнения, а левый~--- за единственность.
\end{lemmanote}

\begin{exercise}
    $A$~--- линейный, то существует обратный ему линейный оператор
\end{exercise}

\begin{exercise}
Пусть $A$~--- ограниченный оператор и у него существует обратный $A^{-1}$. Верно ли, что $A^{-1}$~--- ограниченный?

\textbf{Ответ:} нет, это неверно, пример~--- оператор Вольтера на $C[0;1]$:
\begin{align*}
    (Vf)(x) &=\int_{t=0}^x f(t)dt,\\
    \im \: V &= \{ g \in C^1, g(0)=0 \}.
\end{align*}
Обратный оператор~--- оператор дифференцирования, не является ограниченным
\end{exercise}

\begin{theorem}[8.4, собственная теорема Банаха]\label{self-banah}\label{theorem8.4}
    Пусть $E_1$, $E_2$~--- банаховы пространства,\\$A \in \mathcal{L}(E_1, E_2)$, причём $A$~--- биекция. Тогда $A^{-1}$~--- ограниченный оператор.
\end{theorem}

\begin{proof}
    В общем случае: доказательство через теорему Бэра о категориях, мы доказывать ее не будем.
    Для гильбертовых пространств в параграфе 10 (далее).
\end{proof}

\begin{definition}[Усиление условия $\ker A = \{0\}$]
    Назовем ЛН оператор $A$ \textit{ограниченным снизу}, если
    \[\exists m > 0:\; \forall x \; \|Ax\| \geqslant m \|x\|.\]
\end{definition}

\begin{theorem}[8.1]\label{8.1}
    Пусть $A \in \mathcal{L}(E_1, E_2)$. $\exists A^{-1}$~--- ограниченный $\Leftrightarrow A$~--- ограничен снизу.
\end{theorem}
\begin{proof}
    \begin{itemize}
        \item[$\Leftarrow$]
        Берем $\forall y \in \im\: A,$ для него верно \(\|y\| = \|Ax\| \geqslant m \|x\| = m \|A^{-1}y\|.\)
        Отсюда $\|A^{-1} y\| \leqslant \frac{1}{m} \|y\|$.
        \item[$\Rightarrow$]
        \[ \forall x \in E_1 \; \|x\|=\|A^{-1} y\| \leqslant \| A^{-1}\| \cdot \|y\|=\|A^{-1}\| \cdot \|Ax\| \]
    \end{itemize}
\end{proof}

\begin{exercise}
    $H$~--- гильбертово пространство, $A \in \mathcal{L}(H)$, 
    $\overline{\im A} = H$, $A$~--- ограниченный снизу. Тогда $\im A = H$.
\end{exercise}

\begin{example}[параграф 7, номер 6]
    Доказать, что $\forall g \in C[0;1]$ уравнение $f(x)+\int_0^x f(t)dt=g(x)$ имеет единственное решение, которое непрерывно
    Операторная формулировка: $(I + V)f = g$. По сути нужно найти обратный оператор для $I + V$. ($V$~--- оператор Вольтера)
\end{example}

\begin{example}
    $H, \{e_n\}$~--- ОНБ, $\{f_n\}$~--- ОНС, $\sum \|e_n - f_n\|^2 \le \infty.$ Тогда $\{f_n\}$~--- тоже базис.
\end{example}

Следующие 2 теоремы будут из теории возмущений линейных операторов.

\begin{theorem}[8.2]\label{theorem8.2}
    Пусть $E$~--- банахово пространство, $A \in \mathcal{L}, \|A\| < 1$. Тогда $\exists (I+A)^{-1} \in \mathcal{L}(E)$ 
\end{theorem}
% Вы себя кем считаете, вот вы фивт или как там вас МИ и что-то там ФП еще перед этим. Ну пусть вы фивт.
% <рисует фивт на доске>
\begin{proof}
    Два пути доказательства
    \begin{enumerate}
        \item Упражнение: с помощью теоремы Банаха о сжимающих отображениях. ($x=y-Ax$)
        \item Второй
    \end{enumerate}

    Эвристика: если бы мы работали с числами и необходимо бы было найти $$(1 + A)^{-1} = 1 - A + A^2 + \ldots, \; \|A\| < 1.$$
    Нельзя ли это рассуждение применить для операторов? Рассмотрим последовательность\\$S_n = I - A + A^2 + \ldots + (-1)^nA^n$.
    Так как $\{S_n\}$~--- фундаментальная последовательность:
    \[
    \|S_{n + p} - S_n\| = \left\| \sum_{k = n + 1}^{n + p} (-1)^k A^k\right\| \leq \sum_{k = n + 1}^{n + p} \|A^k\| \leq \sum_{k = n + 1}^{n + p} \|A\|^k \to 0,\; n \to \infty,
    \]
    то последовательность $\{S_n\}$ сходится, так как $\mathcal{L}(E)$~--- банахово (в силу банаховости $E$)

    Утверждается, что $S=R_{I+A}=L_{I+A}$. Рассмотрим \[ (I+A)S = \lim (I+A)S_n = \lim (I+A)(I-A+A^2 +\ldots +(-1)^nA^n)=\lim I+(-1)^n A^{n+1}=I. \]

    Так как $(I+A)S_n = S_n(I+A)$ (многочлены от $A$ перестановочны), то $S$ является как правым так и левым обратным.
    
    Конструкция $(I + A)^{-1} = \sum_{n = 0}^\infty (-1)^n A^n$~--- ряд Неймана.
\end{proof}

\begin{lemmanote}[на вопрос из зала]
    Норма степени оператора неравна степени нормы. Это наблюдается уже на матрицах.
\end{lemmanote}

\begin{lemmanote}[Уточнение теоремы]
Ряд Неймана, появляющийся в доказательстве, сходится, если $\exists n_0 \in \N:\; \|A^{n_0}\| < 1$.
\end{lemmanote}

\begin{example}
    $\|V\|=1$, но $\|V^k\|=\frac 1 {n!}$, поэтому применимо уточнение теоремы
\end{example}

\begin{exercise}
    Оцените сверху нормы операторов 
    $\|(I + A)^{-1}\|$ и $\|I - (I + A)^{-1}\|$.
\end{exercise}

\begin{theorem}[8.3]\label{theorem8.3}
Пусть $E$~--- банахово пространство. \[A, \Delta A \in \mathcal{L}(E),\;\; \exists \; A^{-1} \in \mathcal{L}(E), \|\Delta A\| < \|A^{-1}\|^{-1}.\]

Тогда $\exists (A + \Delta A)^{-1} \in \mathcal{L}(E)$
\end{theorem}

\begin{proof} 
    Рассмотрим $A + \Delta A = A (I + A^{-1} \Delta A)$. Ну и все: у $A$ существует обратный по условию, а у $(I + A^{-1} \Delta A)$ по теореме~\ref{theorem8.2}. Значит и у $A+\Delta A$ существует обратный.
\end{proof}

\begin{exercise}
\begin{enumerate}
    \item 
    \[ \|(A+\Delta A)^{-1}\| \leqslant \frac{\|A^{-1}\|}{1-\|A^{-1}\| \cdot \|\Delta A\|} \]

    \item 
    \[
    \|A^{-1} - (A - \Delta A)^{-1}\| \leqslant 
    \frac{\|A^{-1}\|^2 \cdot \|\Delta A\|}{1 - \|A^{-1}\| \cdot \|\Delta A\|}
    \]
\end{enumerate}
\end{exercise}

\begin{proof}
    \begin{enumerate}
        \item Из конструкции из теоремы \nameref{theorem8.2}: 
        \[ (A+\Delta A)^{-1} = (A(I+A^{-1} \Delta A))^{-1} = (I+A^{-1} \Delta A)^{-1} A^{-1} = \left( \sum_{n=0}^\infty (-1)^n (A^{-1} \Delta A)^n \right) A^{-1}\]

        Пользуясь неравенством треугольника и неравенством $\|AB\| \leq \|A\| \|B\|$ получаем:
        \[ \| (A+ \Delta A)^{-1} \| = \| \left( \sum_{n=0}^\infty (-1)^n (A^{-1} \Delta A)^n \right) A^{-1} \| \leq \| A^{-1} \| \sum_{n=0}^\infty \| A^{-1}\|^n \| \Delta A\|^n = \]

        Так как по условию теоремы $\| A^{-1} \| \| \Delta A \| < 1$, можем применить формулу суммы геометрической прогрессии

        \[ = \frac{\|A^{-1}\|}{1 - \| A^{-1} \| \| \Delta A \|}\]

        \item \[ A^{-1} - (A + \Delta A)^{-1} = (I - (I+A^{-1} \Delta A)^{-1} A^{-1} = (I - \sum_{n=0}^\infty (-1)^n (A^{-1} \Delta A)^n) A^{-1} = (\sum_{n=1}^\infty (-1)^n (A \Delta A)^n) A^{-1} \]

        Дальше аналогично пункту 1 применяя формулу для суммы геометрической прогресси для последовательности без нулевого члена получаем искомое неравенство.
    \end{enumerate}
\end{proof}

\begin{definition}
    Пусть есть $E(\C)$~--- банахово пространство, $A \in \mathcal{L}(E)$, $\lambda \in \C$ называется регулярным значением, если
    \[
    \exists (A - \lambda I)^{-1} \in \mathcal{L}(E).
    \]

    Множетсво регулярных значений называется резольвентным множеством и обозначается $\rho (A)$
\end{definition}

\begin{definition}
    Если $\lambda \in \rho(A)$, то $(A - \lambda I)^{-1}$~--- резольвента.
\end{definition}

% сказка о трех векторах
\begin{definition}
    $\sigma(A)=\C \setminus \rho(A)$~--- спектр.
\end{definition}